\section{Conclusion}
\label{sec:conclusion}

This project reproduced a published KAN-ODE recipe on the Lotka-Volterra
predator-prey system, extended it to two additional dynamical systems (a
damped pendulum and an SIR epidemic model), and diagnosed and fixed two
genuine cross-domain training failures before running five further studies
that each isolated a specific question about solver behavior, gradient
dynamics, basis-function design, stiffness, and comparison against a
fundamentally different modeling paradigm.

The principal quantitative findings are:

\begin{itemize}
\item KAN-ODE extrapolation error on Lotka-Volterra is essentially flat
  across eight limit-cycle periods ($1.1\times$ degradation, training to
  double horizon), while a parameter-matched MLP degrades $8.3\times$ over
  the same extension -- a $40.6\times$ KAN-ODE advantage at the longest
  horizon tested (Table~\ref{tab:extrap-horizon}).
\item SIR's training failure traced to a $19\times$ unit-scale mismatch
  between the network's initial field and the true dynamics, fixable by a
  change of time units and a conservation projection -- two interventions
  that assume nothing about the system beyond what its own equations state.
  Extrapolation $R^2$ moved from $-20.97$ to $+0.9981$.
\item Direct autograd is the correct default for this project's scale:
  adjoint sensitivity already uses less memory at every tested trajectory
  length but costs a consistent $1.9$--$2.1\times$ wall-clock overhead.
\item Gradient-norm roughness does not explain Forward Euler's convergence
  floor (Spearman $\rho=-0.088$, $p=0.87$); the more likely mechanism is a
  systematic bias in the learned field, consistent with the solver ablation's
  own Lipschitz-constant measurements.
\item A learnable blend of B-spline and RBF basis functions never outperforms
  the better pure basis on either system tested.
\item Stiffness-map failures at transitional damping ratios trace to random
  initialization rather than solver order or numerical stiffness, confirmed
  mechanistically by a genuine spurious equilibrium in the learned vector
  field.
\item KAN-ODE is more accurate than sparse symbolic regression at every noise
  level tested, but sparse regression is far more robust ($9\times$ vs.\
  $1540\times$ degradation from clean to $\sigma=0.1$). Two of the SIR
  stability fixes transfer to a real, non-mechanistic epidemic curve; the
  one that transfers best was predicted inapplicable.
\end{itemize}

Every quantitative claim in this report traces to a saved metrics file or a
regenerable figure, and every study that tested a specific mechanistic
hypothesis reports the outcome of a direct control arm rather than the
hypothesis alone.

The validation stage (Section~\ref{sec:phase4}) is the project's principal
open item: the results in Sections~\ref{sec:track-a}--\ref{sec:track-e} rest
on a single random seed and a fixed 10{,}000-epoch training budget.
Completing the epoch-budget sensitivity check and the paper-implementation
audit would extend confidence in every result reported here; the stiffness
study's own experience -- where extending from one seed to three dissolved a
headline finding -- demonstrates concretely why that extension matters.

\FloatBarrier
